% ======================================================================
% 3.1 Равновесие и чувствительность. Ответственный: A — Плешнев Арсений.
%
% Что сдать (main.tex: модель и теория 761–812, порядок 935–982,
% разобранный пример 984–1072; требования к отчёту 1871–1879):
%   1. Проверить размерности членов системы модели; объяснить, почему
%      при mu*c <= lambda равновесия нет.
%   2. Найти q* из условия dot q=0; по закону Литтла выразить T(lambda,mu,c);
%      показать, что T зависит от параметров только через запас s.
%   3. Найти T_lambda, T_c, T_mu и дифференциал dT; вычислить в рабочей точке
%      и пояснить смысл; составить уравнение касательной плоскости T=T(lambda,c).
%   4. Производная по направлению (Delta lambda, Delta c); направление нулевого
%      изменения и доказательство постоянства T вдоль него; уравнения линий
%      уровня и их график в плоскости (lambda,c) при фиксированном mu.
%   5. Градиент T; эластичности E_lambda, E_c; вывод: что сильнее влияет
%      на задержку — рост трафика на 1% или рост ресурса на 1%; знаки.
%   6* (по желанию, на оценку не влияет): худший сценарий методом Лагранжа,
%      delta=0.02; сравнить линейный прогноз с точным значением.
%
% Параметры варианта 3 (приложение Б): mu=80, lambda=600, c=8 (рабочая точка);
% a=1, b=1, p=0.03; (Delta lambda, Delta c)=(20, 0.1); mu2=120, N=16.
% ВНИМАНИЕ: a,b,c приложения А (интегралы) и приложения Б (модель) — разные величины.
%
% Аналитику оформить здесь; численная проверка производных — п. 3.5.1.
% Программные значения для таблицы — task_3_1 (python experiments.py 3.1).
% Рисунок линий уровня: figures/3-1-levels.pdf.
% ======================================================================

\subsection{Равновесие и чувствительность}

\taskhead{Модель и исходные данные}
\todo{A: записать систему уравнений модели и обозначения (main.tex, строки 761–812); подставить рабочую точку варианта 3; вычислить запас $s$ и загрузку $\rho$.}

\taskhead{Размерности и условие равновесия}
\todo{A: проверить размерности членов уравнений модели; объяснить, почему при $\mu c\leq\lambda$ система не может прийти в равновесие.}

\taskhead{Равновесная очередь и задержка}
\todo{A: найти $q^*$ из условия $\dot q=0$; по закону Литтла выразить $T(\lambda,\mu,c)$; показать, что $T$ зависит от параметров только через $s$.}

\taskhead{Производные, дифференциал и касательная плоскость}
\todo{A: вывести $T_\lambda,T_c,T_\mu$ и $dT$; вычислить в рабочей точке и пояснить смысл; составить уравнение касательной плоскости к поверхности $T=T(\lambda,c)$ и показать, что оно задаёт линейный прогноз задержки.}

\taskhead{Производная по направлению и линии уровня}
\todo{A: найти производную по направлению $(\Delta\lambda,\Delta c)=(20,\,0.1)$; найти направление нулевого изменения $T$ и доказать, что вдоль него $T$ постоянна точно; записать уравнения линий уровня.}
\taskfigure[0.75]{figures/3-1-levels.pdf}{Линии уровня задержки в плоскости $(\lambda,c)$ при фиксированном $\mu$; отмечена рабочая точка}{fig:3-1-levels}
\todo{A: пояснить график: несколько положительных значений $T$, подписи линий, рабочая точка.}

\taskhead{Градиент и эластичности}
\todo{A: найти $\operatorname{grad}T$; вычислить $E_\lambda,E_c$; сделать вывод, что сильнее влияет на задержку при одинаковом малом относительном изменении, пояснить знаки эластичностей.}

\taskhead{Сопоставление с программным расчётом}
% Программные значения — вызовы своих функций из solution.py (task_3_1);
% численные производные — из п. 3.5.1. Формат таблицы: main.tex, строки 1905–1915.
\todo{A: листинг вызовов функций и запуск python experiments.py 3.1.}
\begin{table}[htbp]
\centering\footnotesize
\setlength{\tabcolsep}{4pt}
\caption{Сравнение аналитических и программных значений (п. 3.1)}
\label{tab:3-1}
\begin{tabular}{|l|r|r|r|r|}\hline
Величина & \shortstack{Аналитическое\\значение} & \shortstack{Программное\\значение} & \shortstack{Абсолютная\\ошибка} & \shortstack{Относительная\\ошибка}\\\hline
$q^*$ & — & — & — & —\\
$T$ & — & — & — & —\\
$T_\lambda$ & — & — & — & —\\
$T_c$ & — & — & — & —\\
$T_\mu$ & — & — & — & —\\
$E_\lambda$ & — & — & — & —\\
$E_c$ & — & — & — & —\\
\hline\end{tabular}
\end{table}

\taskhead{Вывод}
\todo{A: выводы по чувствительности задержки; расхождения и их причины.}

% Пункт 3.1.6* (худший сценарий, метод Лагранжа, delta=0.02) — по желанию,
% на оценку не влияет. Если выполняется: функция Лагранжа, решение системы,
% обоснование максимума, линейный прогноз и сравнение с точным значением.
